\subsection{Measurable sets and functions}

DEFINITION 7. --- A function f defined on T, with values in a topological space F (Hausdorff or not) is said to be measurable for the measure \(\mu\) (or to be \(\mu\) -measurable) if, for every compact subset K of T, the function \(f_{K}\) is \(\mu_{K}\) -measurable.

This amounts to saying that there exists, for every compact set K, a partition of K into a \(\mu_{K}\) -negligible set N and a sequence \((K_{n})\) of compact sets, such that the restriction of f to each \(K_{n}\) is continuous. Since it is equivalent to say that N is \(\mu_{K}\) -negligible or that it is locally \(\mu\) -negligible (No.~4), one sees that f is \(\mu\) -measurable if and only if, for every compact set K, there exists a partition of K into a locally \(\mu\) -negligible set N and a sequence \((K_{n})\) of compact sets such that \(f_{K_{n}}\) is continuous for all n.~This definition is identical to that of Def. 1 of Ch. IV, §5, No.~1, and one thus recovers the usual concept of measurable function when T is locally compact.

A subset A of T is said to be measurable if its characteristic function is measurable. When A is \(\mu\) -measurable and \(\mu^{\bullet}(A) < +\infty\) , this number is denoted simply \(\mu(A)\) and is called the measure of A. One similarly writes \(\mu(f)\) for \(\mu^{\bullet}(f)\) when f is \(\geqslant 0\) , \(\mu\) -measurable, and \(\mu^{\bullet}(f) < +\infty\) .

If \(\theta\) is a complex measure on \(\mathbf{T}\) , a function \(f\) (resp. a subset of \(\mathbf{T}\) ) is said to be \(\theta\) -measurable if it is measurable for the positive measure \(|\theta|\) . The results below may be extended to complex measures.

Example. --- Let L be a compact subset of T, \(\lambda\) a measure on L, and \(\mu\) the measure on T defined by \(\lambda\) (No.~3, Example 2). A function f defined on T is \(\mu\) -measurable if and only if \(f_{L}\) is \(\lambda\) -measurable. For, this condition is obviously necessary. Conversely, if it is satisfied, there exists a partition of L into a \(\lambda\) -negligible set N and a sequence \((L_{n})\) of compact sets, such that \(f_{L_{n}}\) is continuous for all n.~If K is a compact subset of T, the set \(K - \bigcup_{n}(K \cap L_{n})\) has an
intersection with L that is \(\lambda\) -negligible, hence this set is \(\mu\) -negligible by formula (3) of No.~3, and the restriction of f to \(K \cap L_{n}\) is continuous for all n.

Def. 7 permits extending, without a new proof, a number of results on measurable functions to the case of spaces not locally compact. Here are some of them, which we shall use henceforth without further reference: the open sets and the closed sets of T are \(\mu\) -measurable; the \(\mu\) -measurable sets form a tribe (Ch. IV, §5, No.~4, Cor. 2 of Th. 2), that contains the Borel sets of T (loc. cit., Cor. 3), and the Souslin sets (Ch. IV, §5, No.~1, Cor. 2 of Prop. 3) \(^{(1)}\) . The usual algebraic operations on numerical functions preserve measurability (Ch. IV, §5, No.~3), as do the operations of countable passage to the limit (loc. cit., No.~4, Th. 2 and Cor. 1). The following property merits more explicit mention:

PROPOSITION 4. --- Let \(f\) be a positive function and \((g_n)_{n \geqslant 1}\) a sequence of \(\mu\) -measurable positive functions on \(T\) . Setting \(g = \sum_{n \geqslant 1} g_n\) , one has

\[
\mu^ {\bullet} (f g) = \sum_ {n \geqslant 1} \mu^ {\bullet} (f g _ {n}).\tag{4}
\]

Set \(h_n = \sum_{i=1}^{n} g_i\) for all \(n \geqslant 1\) . For every compact subset \(K\) of \(T\) ,

\[
\mu_ {\mathrm{K}} ^ {\bullet} \big ((f h _ {n}) _ {\mathrm{K}} \big) = \sum_ {i = 1} ^ {n} \mu_ {\mathrm{K}} ^ {\bullet} \big ((f g _ {i}) _ {\mathrm{K}} \big)
\]

by Prop. 2 of Ch. V, §1, No.~1 applied to the compact space K. Passing to the limit with respect to the increasing directed set of compact subsets of T, one obtains

\[
\mu^ {\bullet} (f h _ {n}) = \sum_ {i = 1} ^ {n} \mu^ {\bullet} (f g _ {i}).
\]

Now, \(fg\) is the limit of the increasing sequence \((fh_n)_{n\geqslant 1}\) , whence \(\mu^{\bullet}(fg) = \lim_{n\to \infty}\mu^{\bullet}(fh_n)\) ; the preceding formula then immediately implies (4).

COROLLARY. --- Let \((\mathbf{A}_{n})\) be a sequence of pairwise disjoint measurable subsets, with union A. For every subset B of T,

\[
\mu^ {\bullet} (\mathrm{A} \cap \mathrm{B}) = \sum_ {n} \mu^ {\bullet} (\mathrm{A} _ {n} \cap \mathrm{B})
\]

and in particular

\[
\mu^ {\bullet} (\mathrm{A}) = \sum_ {n} \mu^ {\bullet} (\mathrm{A} _ {n}).
\]

Among the properties of measurable functions or sets that extend as above to Hausdorff spaces, we cite also Prop. 12 of Ch. IV, §5, No.~8 ( \(\mu\) -dense families of compact sets). Thus, a function f with values in a topological space (Hausdorff or not) is \(\mu\) -measurable if and only if the set of compact subsets K of T, such that \(f_{K}\) is continuous, is \(\mu\) -dense (loc. cit., No.~10, Prop. 15).

